\subsection{Cross-Scenario Comparative Analysis}

This section synthesizes the empirical data collected across the four experimental phases. By evaluating the smart home network's operational baseline against a set of different attacks scenarios, we can contrast the structural and behavioural footprint of each attack vector. The comparative analysis highlights a critical reality in IoT network forensics. \textbf{The severity of a network compromise does not strictly correlate with flow volume}. Adversaries adapt their techniques and procedures, ranging from massive volumetric spikes to highly evasive port masquerading, based on their specific operational objectives.

\subsubsection{The Macro-Statistics and Detection Difficulty Matrix}

To establish a view of the threat landscape, the primary network metrics for all four scenarios have been aggregated. Table \ref{tab:cross_scenario_macro} illustrates how standard flow analysis succeeds against noisy attacks but fails against advanced stealth tactics.

\begin{table}[H]
	\centering
	\renewcommand{\arraystretch}{1.3}
	\resizebox{\textwidth}{!}{%
		\begin{tabular}{@{}l l r r l@{}}
			\toprule
			\textbf{Scenario} & \textbf{Threat Profile} & \textbf{Total Flows} & \textbf{Volumetric $\Delta$} & \textbf{Primary Forensic Indicator} \\ 
			\midrule
			\textbf{Scenario 1} & Operational Baseline & 181 & 0.0\% & 100\% adherence to MUD profiles. \\
			\textbf{Scenario 2} & Recon \& Brute-Force & 5,987 & +3,329\% & Massive global flow density surge. \\
			\textbf{Scenario 3} & Stealth Exfiltration & 199 & +10.5\% & DNS anomalies and Port 443 Masquerading. \\
			\textbf{Scenario 4} & Botnet Infection & 210 & +130.8\% & Internal lateral scanning and outbound UDP flood. \\
			\bottomrule
		\end{tabular}%
	}
	\caption{Macro-Statistics and Detection Difficulty Matrix across all experimental scenarios.}
	\label{tab:cross_scenario_macro}
\end{table}

\subsubsection{Volumetric Visibility vs. Architectural Stealth}

The most striking contrast in the experimental data is the volumetric disparity between the different attack methodologies. 

Both Scenario 2 and Scenario 4 performed an automated subnet scan to map the environment. Because network scanners are noisy by flooding targets with thousands of simultaneous \texttt{SYN} probes, this attack is highly visible and easily caught by rudimentary threshold-based intrusion detection systems.

Conversely, Scenario 3 demonstrates the danger of sophisticated post-exploitation tactics designed to bypass volumetric heuristics. By relying on a single targeted payload to trigger the malware, the attacker generated a negligible 10.5\% increase in total flows (199 connections). Rather than creating new network noise, the compromised Smart Camera simply redirected its existing high-bandwidth UDP video stream to an unauthorized external IP. Because the aggregate network throughput remained seemingly stable, this scenario proves that deep packet metadata inspection—specifically looking for port masquerading and hijacking DNS polling frequencies—is strictly required to detect advanced exfiltration.

\subsubsection{The TCP Connection State Tripwire}

The distribution of TCP connection states proved to be one of the most reliable automated tripwires for detecting lateral movement within the zero-trust architecture.

In the secure Scenario 1 baseline, the environment was heavily dominated by the \texttt{SF} (Normal) state flag, indicating successful connections and clean teardowns. However, when an adversary attempts to breach internal boundaries, this ratio violently inverts:

\begin{itemize}
	\item \textbf{External Lateral Scanning (Scenario 2):} As the attacker node flooded the IoT edge devices with unauthorized probes, the devices actively defended themselves by issuing TCP Reset (RST) packets, resulting in an explosive 6,041 flows terminating in a \texttt{REJ} (Rejected) state.
	\item \textbf{Internal Insider Threat (Scenario 4):} A similar structural inversion occurred when the Smart Thermostat was conscripted into a botnet. As the hijacked thermostat aggressively scanned the local \texttt{10.0.0.0/24} subnet, Zeek recorded 128 malicious flows terminating in a \texttt{REJ} status.
\end{itemize}

Crucially, Scenario 3 entirely bypassed this tripwire. Because the attacker utilized a single payload and did not indiscriminately probe closed ports, the network generated zero \texttt{REJ} states, maintaining an illusion of a mathematically healthy TCP footprint with exactly 132 \texttt{SF} flows.

\subsubsection{The Degradation of Zero-Trust and MUD Profiles}

The progression from Scenario 1 to Scenario 4 effectively maps the degradation of the Manufacturer Usage Description (MUD) profiles under increasing adversarial pressure. 

The control environment established a strict node-to-node isolation matrix. Scenario 2 violated this via unauthorized probes directed at the devices. Scenario 3 bypassed these rules via egress manipulation, utilizing DNS hijacking and routing stolen video feeds to an unauthorized external rogue cloud. Finally, Scenario 4 represented a total collapse of both ingress and egress trust boundaries, as the Smart Thermostat established an unauthorized IRC-style Command and Control link and weaponized its throughput to participate in a global DDoS campaign.